\section{Bài toán chung và khung phân tích}

\subsection{Bài toán chung}

Bài toán chung của nhóm là:

\begin{quote}
\textit{Trong giai đoạn 2000--2023, mức độ phát triển của các quốc gia và khu vực trên thế giới thay đổi như thế nào khi xét đồng thời các khía cạnh kinh tế, giáo dục, sức khỏe--dân số và môi trường? Các khác biệt giữa khu vực và nhóm thu nhập cho thấy những cơ hội hoặc thách thức nào đối với phát triển bền vững?}
\end{quote}

Bài toán này yêu cầu nhóm không chỉ nhìn vào tăng trưởng kinh tế, mà còn kiểm tra liệu sự phát triển đó có đi kèm với cải thiện về giáo dục, sức khỏe và điều kiện môi trường hay không. Do đó, mỗi thành viên phụ trách một hướng phân tích riêng, nhưng các kết quả cuối cùng sẽ được tổng hợp để trả lời cùng một câu hỏi trung tâm.

\subsection{Các câu hỏi phân tích thành phần}

\begin{longtable}{L{4cm}L{10cm}}
\caption{Vai trò của bốn hướng phân tích trong bài toán chung. Các hướng phân tích được thiết kế để bổ sung cho nhau thay vì trả lời bốn câu hỏi độc lập.}
\label{tab:analysis-framework}\\
\toprule
Hướng phân tích & Vai trò trong bài toán chung \\
\midrule
\endfirsthead
\toprule
Hướng phân tích & Vai trò trong bài toán chung \\
\midrule
\endhead
Phát triển kinh tế & Đánh giá nền tảng vật chất của phát triển thông qua GDP, GDP bình quân đầu người, tăng trưởng GDP và quan hệ với tỷ lệ nghèo. \\
Giáo dục & Kiểm tra khả năng tiếp cận giáo dục và đầu tư cho giáo dục, từ đó đánh giá chất lượng nguồn nhân lực trong dài hạn. \\
Sức khỏe và dân số & Đánh giá kết quả phát triển về mặt đời sống con người thông qua tuổi thọ, tử vong trẻ em, chi tiêu y tế và quy mô dân số. \\
Môi trường và biến đổi khí hậu & Kiểm tra mặt trái hoặc áp lực của phát triển thông qua phát thải CO\textsubscript{2}, năng lượng tái tạo và diện tích rừng. \\
\bottomrule
\end{longtable}

\subsection{Logic kết nối các phần}

Kinh tế cho biết một quốc gia có tăng trưởng và tạo ra nguồn lực hay không. Giáo dục và sức khỏe cho biết nguồn lực đó có gắn với cải thiện chất lượng con người hay không. Môi trường cho biết quá trình phát triển có tạo ra áp lực lên tài nguyên và khí hậu hay không. Khi đặt bốn nhóm chỉ số cạnh nhau trong Tableau, nhóm có thể so sánh các quốc gia và khu vực theo hướng tổng quan hơn: phát triển nhanh, phát triển bao trùm, phát triển cải thiện đời sống, và phát triển có tính bền vững. Vì dữ liệu WDI là dữ liệu quan sát ở cấp quốc gia, báo cáo chỉ diễn giải các mẫu hình và tương quan mô tả, không xem các biểu đồ là bằng chứng nhân quả.
